
\section{Conclusion}

MTI improves employment opportunities if it is complemented with favorable institutions and policies. In the absence of favorable social infrastructure, MTI does not increase employment outcomes even when it increases cognitive ability.  

MTI has increased salary employment in Ethiopia and this is an expected outcome of the key feature of the reconfiguration of the Ethiopian state along ethnic lines. The country's 1994 constitution created local states with their own official ethnic languages, increasing the prospects of securing salaried positions by those that obtain their education in their native tongues, since the various functions of the state governments are performed using local languages. 

However, MTI has not performed well in terms of wage employment, most likely because the overall social infrastructure in Ethiopia does not support a more robust implementation of the MTI policy. Employees do not have legal rights to work in their native tongues in the private sector, and wage employers have continued to use Amharic in conducting their day-to-day operations, restricting the employment opportunities of those who are educated in their mother tongues and may have lost some fluency in Amharic as a result. 

The impact of MTI on Ethiopia’s development can be optimized by further devolution of power from the political center to the regional states. For instance, it may be appropriate for local governments to look into instituting labor laws that afford workers the legal rights to work in their native languages, similar to what other multilingual countries (e.g., Canada, Switzerland, etc.) have successfully implemented. 

These laws and policies can be instituted differently in different states based on the needs and capacities of each state. It would be prudent to weigh the benefits of these measures against their potential adverse effect on economic efficiency, which may result from undue restrictions on labor mobility that new laws and policies may impose. Where there are constraints on well-trained local talent, regional states can adopt laws that provide exceptions to employers to allow them to attract the requisite talent from the national labor market. Therefore, while aiming to reduce or eliminate the loss of productivity embedded in the status quo due to a preexisting bias in favor of Amharic, which, as we have shown in this paper, denies more qualified local talent from fully participating in the labor market, policy makers in the state governments may need to proceed cautiously to avoid instituting measures that could lead to inferior economic outcomes.
